% GENERATED FILE. Edit canonical lessons, then run npm run generate:books.
% canonical-source-hash: fnv1a64:666e74b0df197e5b
% canonical-lessons: SA-C228-itihasah, SA-C228-bhugolah, SA-C228-sangitam, SA-C228-nrtyam, SA-C228-citrakala

\chapter{History, Geography, Music, Dance, Painting}
\label{ch:sa-history-geography-music-dance-painting}
\hlchaptermodality{228}

\begin{chapteropening}
\textbf{By the end of this chapter:} \emph{I can say history, geography, music, dance and painting.}

Complete the last lesson of chapter 228: I can say history, geography, music, dance and painting.
\end{chapteropening}

\section[\texorpdfstring{\sk{इतिहासः} (itihāsaḥ)}{इतिहासः (itihāsaḥ)}]{\sk{इतिहासः} (itihāsaḥ) — history}

\label{lesson:SA-C228-itihasah}



\begin{quote}
Before the new one: say the Sanskrit for a beak, then the Sanskrit for a vulture.
\end{quote}

\subsection*{You'll want to know: \sk{इतिहासः}}
\textbf{\sk{इतिहासः}} — \emph{itihāsaḥ} — \textquotedblleft{}history\textquotedblright{}.

\begin{grammarlens}[title={What you've built}]
One more word for this chapter.
\end{grammarlens}

\subsection*{Your turn}
\begin{itemize}
\raggedright
  \item \emph{Say it:} \emph{itihāsaḥ}
  \item \emph{Say it:} \emph{itihāsaḥ}, once more
  \item \emph{Say it:} say \emph{itihāsaḥ}
  \item \emph{Recall:} say the Sanskrit for a beak, then the Sanskrit for a vulture, then say \emph{itihāsaḥ} again
\end{itemize}

\begin{tcolorbox}[breakable,colback=teal!4,colframe=teal!35!black,title={Before you move on}]
What does \sk{इतिहासः} mean? (\textquotedblleft{}History\textquotedblright{}.) Say it once more.
\end{tcolorbox}
\section[\texorpdfstring{\sk{भूगोलः} (bhūgolaḥ)}{भूगोलः (bhūgolaḥ)}]{\sk{भूगोलः} (bhūgolaḥ) — geography}

\label{lesson:SA-C228-bhugolah}



\begin{quote}
Before the new one: say the Sanskrit for a vulture, then the Sanskrit for history.
\end{quote}

\subsection*{You'll want to know: \sk{भूगोलः}}
\textbf{\sk{भूगोलः}} — \emph{bhūgolaḥ} — \textquotedblleft{}geography\textquotedblright{}.

\begin{grammarlens}[title={What you've built}]
One more word for this chapter.
\end{grammarlens}

\subsection*{Your turn}
\begin{itemize}
\raggedright
  \item \emph{Say it:} \emph{bhūgolaḥ}
  \item \emph{Say it:} \emph{bhūgolaḥ}, once more
  \item \emph{Say it:} say \emph{bhūgolaḥ}
  \item \emph{Recall:} say the Sanskrit for a vulture, then the Sanskrit for history, then say \emph{bhūgolaḥ} again
\end{itemize}

\begin{tcolorbox}[breakable,colback=teal!4,colframe=teal!35!black,title={Before you move on}]
What does \sk{भूगोलः} mean? (\textquotedblleft{}Geography\textquotedblright{}.) Say it once more.
\end{tcolorbox}
\section[\texorpdfstring{\sk{संगीतम्} (saṅgītam)}{संगीतम् (saṅgītam)}]{\sk{संगीतम्} (saṅgītam) — music}

\label{lesson:SA-C228-sangitam}



\begin{quote}
Before the new one: say the Sanskrit for history, then the Sanskrit for geography.
\end{quote}

\subsection*{You'll want to know: \sk{संगीतम्}}
\textbf{\sk{संगीतम्}} — \emph{saṅgītam} — \textquotedblleft{}music\textquotedblright{}.

\begin{grammarlens}[title={What you've built}]
One more word for this chapter.
\end{grammarlens}

\subsection*{Your turn}
\begin{itemize}
\raggedright
  \item \emph{Say it:} \emph{saṅgītam}
  \item \emph{Say it:} \emph{saṅgītam}, once more
  \item \emph{Say it:} say \emph{saṅgītam}
  \item \emph{Recall:} say the Sanskrit for history, then the Sanskrit for geography, then say \emph{saṅgītam} again
\end{itemize}

\begin{tcolorbox}[breakable,colback=teal!4,colframe=teal!35!black,title={Before you move on}]
What does \sk{संगीतम्} mean? (\textquotedblleft{}Music\textquotedblright{}.) Say it once more.
\end{tcolorbox}
\section[\texorpdfstring{\sk{नृत्यम्} (nṛtyam)}{नृत्यम् (nṛtyam)}]{\sk{नृत्यम्} (nṛtyam) — dance}

\label{lesson:SA-C228-nrtyam}



\begin{quote}
Before the new one: say the Sanskrit for geography, then the Sanskrit for music.
\end{quote}

\subsection*{You'll want to know: \sk{नृत्यम्}}
\textbf{\sk{नृत्यम्}} — \emph{nṛtyam} — \textquotedblleft{}dance\textquotedblright{}.

\begin{grammarlens}[title={What you've built}]
One more word for this chapter.
\end{grammarlens}

\subsection*{Your turn}
\begin{itemize}
\raggedright
  \item \emph{Say it:} \emph{nṛtyam}
  \item \emph{Say it:} \emph{nṛtyam}, once more
  \item \emph{Say it:} say \emph{nṛtyam}
  \item \emph{Recall:} say the Sanskrit for geography, then the Sanskrit for music, then say \emph{nṛtyam} again
\end{itemize}

\begin{tcolorbox}[breakable,colback=teal!4,colframe=teal!35!black,title={Before you move on}]
What does \sk{नृत्यम्} mean? (\textquotedblleft{}Dance\textquotedblright{}.) Say it once more.
\end{tcolorbox}
\section[\texorpdfstring{\sk{चित्रकला} (citrakalā)}{चित्रकला (citrakalā)}]{\sk{चित्रकला} (citrakalā) — painting}

\label{lesson:SA-C228-citrakala}



\begin{quote}
Before the new one: say the Sanskrit for music, then the Sanskrit for dance.
\end{quote}

\subsection*{You'll want to know: \sk{चित्रकला}}
\textbf{\sk{चित्रकला}} — \emph{citrakalā} — \textquotedblleft{}painting\textquotedblright{}.

\begin{grammarlens}[title={What you've built}]
One more word for this chapter.
\end{grammarlens}

\subsection*{Your turn}
\begin{itemize}
\raggedright
  \item \emph{Say it:} \emph{citrakalā}
  \item \emph{Say it:} \emph{citrakalā}, once more
  \item \emph{Say it:} say \emph{citrakalā}
  \item \emph{Recall:} say the Sanskrit for music, then the Sanskrit for dance, then say \emph{citrakalā} again
\end{itemize}

\begin{tcolorbox}[breakable,colback=teal!4,colframe=teal!35!black,title={Before you move on}]
What does \sk{चित्रकला} mean? (\textquotedblleft{}Painting\textquotedblright{}.) Say it once more.
\end{tcolorbox}
